% 第 54 章骨架：导引框已定，正文待写。写作计划见 docs/PLAN.md“第 54 章”，写法见 docs/STYLE.md。
\chapter{新凯恩斯：名义刚性}\label{ch:newkeynesian}

\begin{chapterguide}
\guidequestion{理性的企业为什么不随时调价？}
\begin{guidegoals}
\item 由泰勒（1980）的交错合同与卡尔沃（1983）定价推出新凯恩斯菲利普斯曲线；
\item 解释菜单成本（Mankiw 1985）与近似理性（Akerlof \& Yellen 1985）为什么能产生大的总量效应；
\item 区分名义刚性与实际刚性；
\item 用微观价格数据检验调价频率。
\end{guidegoals}
\guideprereq{\chref{ch:rbc}\enspace\nameref{ch:rbc}}
\end{chapterguide}

\chapterstub
